

Your colleague Cynthia likes writing asides in messages (but then again,
    don't we all (and there's nothing strictly wrong with that anyway)). As
some people find the text hard to read this way, we are going to try to
extract some parts of the sentences out into equally-sized footnotes instead
of brackets.

For example, \texttt{this(complicated(nested(passage)))} could become:

\indent \texttt{this[1]} \\
\indent \texttt{[1] complicated[2]} \\
\indent \texttt{[2] nested(passage)}

Given one such message, an unbroken string of up to $10^6$ characters
(letters and brackets, without any spaces), isolate exactly
two-thirds of the Latin letters into footnotes, each footnote containing
exactly the same number of Latin letters.

